% ===========================================================================
% The LeTrocq blueprint.
%
%   Chapter 1  = DESIGN.md      — the graded parametricity translation
%   Chapter 2  = bidir_solver.md — the demand-driven, witness-producing pass
%
% Every environment is tied to its Lean declaration with \lean{...}; \leanok
% marks it as formalised (all of LeTrocq is, so the graph is fully green) and
% \uses{...} draws the dependency edges. Fully-qualified names live in the
% `LeTrocq` namespace (`LeTrocq.Param`, `LeTrocq.Transfer.assemble`, ...).
% ===========================================================================

\chapter{The graded parametricity translation}
\label{chap:translation}

LeTrocq is a Lean-native, \emph{graded} reformulation of the parametricity /
proof-transfer translation of Trocq (Cohen--Crance--Mahboubi), in the
\emph{no-univalence} fragment. The translation is \textbf{twofold}: a
counterpart translation $\Cpt{\cdot}$ and a relational translation
$\Rel{\cdot}$, tied by the abstraction theorem
\[
  t : A \quad\Longrightarrow\quad \Rel{t} : \Sem{A}\, t\, \Cpt{t},
  \qquad\text{where } \Sem{A} := \Rel{A}.R .
\]
$\Cpt{t}$ rebuilds $t$ on the target side; $\Rel{t}$ witnesses that $\Cpt{t}$
really is the counterpart. Neither ever unfolds a definition --- both bottom out
at \texttt{@[trocq]}-registered witnesses (\Cref{sec:leaves}).

\section{The classical (ungraded) translation}
\label{sec:classical}

Ignoring grading, on the Calculus of Constructions the two translations are the
standard binary parametricity translation.

\begin{definition}[The ungraded translation]
  \label{def:classical}
  The counterpart $\Cpt{\cdot}$ and relation $\Rel{\cdot}$ are defined by
  mutual structural recursion:
  \[
    \begin{array}{llcll}
      \Cpt{x}          &= x'                 &\qquad& \Rel{x}         &= x_R \\
      \Cpt{t\,u}       &= \Cpt{t}\,\Cpt{u}   &&      \Rel{t\,u}      &= \Rel{t}\; u\; \Cpt{u}\; \Rel{u} \\
      \Cpt{\lambda x{:}A.\,t} &= \lambda x'{:}\Cpt{A}.\,\Cpt{t} && \Rel{\lambda x{:}A.\,t} &= \lambda x\,x'\,x_R.\,\Rel{t} \\
      \Cpt{\Pi x{:}A.\,B}     &= \Pi x'{:}\Cpt{A}.\,\Cpt{B}     && \Rel{\Pi x{:}A.\,B} &= \Pi x\,x'\,x_R.\,\Rel{B} \\
      \Cpt{\square}    &= \square            &&      \Rel{\square}   &= \square \to \square \to \square
    \end{array}
  \]
  The application clause $\Rel{t\,u} = \Rel{t}\,u\,\Cpt{u}\,\Rel{u}$ is the
  \emph{abstraction theorem}, and drives every structural case below. Here
  $\Rel{A}$ is a bare relation $A \to \Cpt{A} \to \square$; the grading of
  \Cref{sec:hierarchy} refines it into a record.
\end{definition}

\section{The graded hierarchy}
\label{sec:hierarchy}

To carry transport \emph{maps} (not just a relation), $\Rel{\cdot}$ on a type
returns a record --- a witness of \Cref{def:param} --- rather than a bare
relation.

\begin{definition}[The map-class diamond]
  \label{def:map-class}
  \lean{LeTrocq.MapClass}
  \leanok
  $\MapClass$ is the finite order $0 < 1 < \{2a, 2b\} < 3 < 4$, with a
  decidable order \texttt{le}, a join and a meet. Semantically a class records
  how much structure a relation carries: bare relation ($0$), a forward map
  ($1$), soundness \texttt{map\_in\_R} ($2a$), completeness \texttt{R\_in\_map}
  ($2b$), both ($3$), and their mutual-inverse coherence ($4$).
\end{definition}

\begin{definition}[Parametricity classes]
  \label{def:param-class}
  \lean{LeTrocq.ParamClass}
  \leanok
  \uses{def:map-class}
  A $\mathsf{ParamClass}$ is a \emph{pair} $\cls{m}{n}$ of map-classes, one per
  transport direction, with pointwise order/join/meet, a direction swap
  \texttt{negate}, and a predicate \texttt{requiresAxiom} that holds iff some
  component is $\geq 2b$ (i.e.\ the class needs univalence on a sort or funext on
  a $\Pi$).
\end{definition}

\begin{definition}[The six graded records]
  \label{def:map-has}
  \lean{LeTrocq.MapHas}
  \leanok
  \uses{def:map-class}
  For a relation $R : A \to B \to \Type$, $\MapHas\;c\;R$ is the structure a
  relation carries at class $c$, selecting one of the records
  \texttt{Map0Has} $\ldots$ \texttt{Map4Has}: \texttt{map} $:A\to B$
  (from $1$), \texttt{map\_in\_R} (from $2a$), \texttt{R\_in\_map} (from $2b$),
  and \texttt{R\_in\_mapK} (the coherence, at $4$). All six live at the uniform
  \SortU\ level $\max(u, v{+}1)$, so a \PropSort\ object fits with no \texttt{ULift}.
\end{definition}

\begin{definition}[The witness record]
  \label{def:param}
  \lean{LeTrocq.Param}
  \leanok
  \uses{def:map-has, def:param-class}
  $\Param\;m\;n\;A\;B$ bundles a relation $R : A \to B \to \Type$ with
  $\texttt{cov} : \MapHas\;m\;R$ (forward structure) and
  $\texttt{contra} : \MapHas\;n\;(\mathrm{flip}\;R)$ (backward structure). Thus
  $\cls{4}{4}$ is a full type equivalence, $\cls{1}{0}$ a forward function,
  $\cls{0}{1}$ a backward one, $\cls{0}{0}$ a bare relation. Symmetry
  (\texttt{Param.sym}) swaps \texttt{cov}/\texttt{contra} onto the reversed
  relation.
\end{definition}

\begin{lemma}[Class-4 relations are subsingletons]
  \label{lem:map4-subsingleton}
  \lean{LeTrocq.Map4Has.subsingleton}
  \leanok
  \uses{def:map-has}
  If $R$ carries \texttt{Map4Has} data, then $R\,a\,b$ is a subsingleton for all
  $a, b$.
\end{lemma}
\begin{proof}
  \leanok
  Two inhabitants of $R\,a\,b$ both map, via \texttt{R\_in\_map}, to proofs of
  $\texttt{map}\,a = b$, which are equal by Lean's proof irrelevance; the
  coherence \texttt{R\_in\_mapK} then forces the two inhabitants equal. This is
  exactly ``no univalence $\Rightarrow$ class $4$ $=$ class $3$ on h-props'',
  and it makes the $\cls{4}{4}$ coherence \emph{free} on any relation reachable
  from class-4 data.
\end{proof}

\begin{definition}[Weakening]
  \label{def:weaken}
  \lean{LeTrocq.Param.weaken}
  \leanok
  \uses{def:param, def:map-class}
  The diamond is generated from six covering-edge forgets
  ($4 \to 3 \to 2a \to 1 \to 0$ and $3 \to 2b \to 1$), each a bare record
  projection. Composing them, $\texttt{MapClass.weaken}\;s\;t$ maps
  $\MapHas\;s\;R \to \MapHas\;t\;R$ whenever $t \leq s$ (impossible cases
  discharged by \texttt{nomatch}); \texttt{Param.weaken} lifts this
  componentwise to $\Param\;s_m\;s_n \to \Param\;t_m\;t_n$. This is the bridge of
  the ``register strong, use minimal'' discipline (\Cref{sec:grading}).
\end{definition}

\section{The relation and the universe crux}
\label{sec:universe}

\begin{definition}[The carried relation]
  \label{def:sem}
  \uses{def:param}
  $\Sem{A} := \Rel{A}.R$: the relation projected off the graded witness the
  translation builds for a type. It is all a term position ever consumes of a
  type, and it is grade-invariant.
\end{definition}

The universe clause is the crux. On a type, $\Rel{\cdot}$ returns a record, and
the relation carried by the universe is $\Param$ \emph{itself}:
$\Sem{\square} = \lambda A\,A'.\;\Param\;\cls{p}{q}\;A\;A'$, so
$\Rel{A} : \Sem{\square}\,A\,\Cpt{A} = \Param\;\cls{p}{q}\;A\;\Cpt{A}$ --- the
type witness $\Rel{A}$ is literally an element of the universe's own relation,
one level up. The inner class $\cls{p}{q}$ records how strongly the bound type
variable is related; it is \emph{independent} of the (capped) outer class.

\begin{definition}[The reflexive witness]
  \label{def:param-refl}
  \lean{LeTrocq.paramRefl}
  \leanok
  \uses{def:param}
  $\texttt{paramRefl}\;A : \Param\;\cls{4}{4}\;A\;A$ with relation
  $R\,a\,b := \mathsf{PLift}\,(a = b)$: the identity map is an equivalence, so all
  four fields hold (coherence by \texttt{Eq} casing). Weakening it gives
  $\texttt{paramIdAt}\;p\;q$ at any class, choice-free.
\end{definition}

\begin{theorem}[The universe combinator, capped at $\cls{2a}{2a}$]
  \label{thm:param-type}
  \lean{LeTrocq.paramTypeAtInner}
  \leanok
  \uses{def:param-refl, def:weaken, def:param-class}
  $\texttt{paramTypeInner}\;p\;q : \Param\;\cls{2a}{2a}\;\Type_w\;\Type_w$ has
  relation $\lambda A\,A'.\;\Param\;\cls{p}{q}\;A\;A'$, with \texttt{map\_in\_R}
  the $\texttt{Eq.rec}$ of $\texttt{paramIdAt}\;p\;q$. It cannot climb above
  $\cls{2a}{2a}$: the completeness field would demand \emph{univalence}. The
  inner class $\cls{p}{q}$ is free and independent of the cap.
  $\texttt{paramTypeAtInner}\;m\;n\;p\;q$ weakens it to any outer class
  $\leq \cls{2a}{2a}$.
\end{theorem}

\begin{theorem}[The \PropSort\ universe, at $\cls{4}{4}$]
  \label{thm:param-prop}
  \lean{LeTrocq.paramProp}
  \leanok
  \uses{def:param}
  $\texttt{paramProp} : \Param\;\cls{4}{4}\;\PropSort\;\PropSort$ with relation
  $\lambda P\,P'.\;\mathsf{PLift}\,(P \leftrightarrow P')$. Unlike $\Type$, it
  reaches the full equivalence: \texttt{map\_in\_R} is $\texttt{Eq.rec}$,
  completeness \texttt{R\_in\_map} is \textbf{\texttt{propext}}, and the coherence
  is free by proof irrelevance (the relation is a subsingleton).
\end{theorem}

\begin{definition}[Propositions relate by equivalence]
  \label{def:param-of-iff}
  \lean{LeTrocq.paramOfIff}
  \leanok
  \uses{thm:param-prop}
  $\texttt{paramOfIff} : (P \leftrightarrow P') \to \Param\;\cls{4}{4}\;P\;P'$
  carries the equivalence as its relation $\mathsf{PLift}\,(P \leftrightarrow P')$;
  completeness is proof irrelevance, so it needs no axiom beyond the given
  \texttt{iff}. Its inverse on maps is \texttt{iffOfParam}. This is the shared
  builder every \PropSort-valued relator (\texttt{Eq}, the connectives) hands its
  result to.
\end{definition}

\section{The structural combinators}
\label{sec:combinators}

Each combinator builds a $\Param$ for a former out of $\Param$s for its parts,
at any output class, with the parts at the class a grading table dictates
(\Cref{sec:grading}).

\begin{definition}[The arrow relation and its grading table]
  \label{def:arrow-variance}
  \lean{LeTrocq.arrowVariance}
  \leanok
  \uses{def:param-class}
  $\RArrow\;R_A\;R_B\;f\;f' := \forall a\,a'.\,R_A\,a\,a' \to R_B\,(f a)\,(f' a')$
  is the respectful relation. $\texttt{arrowVariance}\;c$ returns the minimal
  (domain, codomain) classes needed to build the arrow at output class $c$,
  computed from the per-map-class table \texttt{mapArrowVariance} (verbatim from
  Trocq's \texttt{class.elpi}) as $\mathrm{join}$ of the cov requirement with the
  negated contra one.
\end{definition}

\begin{theorem}[The graded arrow]
  \label{thm:param-arrow}
  \lean{LeTrocq.paramArrow}
  \leanok
  \uses{def:arrow-variance, def:param, def:weaken, lem:map4-subsingleton}
  $\texttt{paramArrow}\;m\;n$ builds $\Param\;\cls{m}{n}\;(A \to B)\;(A' \to B')$
  from parts at the \texttt{arrowVariance}-minimal classes, for \emph{every}
  output class including $\cls{4}{4}$. The forward map is
  $B.\mathrm{fwd} \circ f \circ A.\mathrm{bwd}$; completeness in each direction is
  a \texttt{funext} of the parts' completeness; the $\cls{4}{4}$ coherence is free
  because class-4 parts have subsingleton relations
  (\Cref{lem:map4-subsingleton}), so the arrow relation is a subsingleton.
\end{theorem}

\begin{definition}[The dependent $\Pi$ relation and its grading table]
  \label{def:forall-variance}
  \lean{LeTrocq.forallVariance}
  \leanok
  \uses{def:param-class}
  With a codomain family $B : A \to \SortU$ and an indexed relation
  $R_B : \forall a\,a'.\,R_A\,a\,a' \to B\,a \to B'\,a' \to \Type$, the
  $\Pi$-relation is
  $\RForall\;R_A\;R_B\;f\;f' := \forall a\,a'\,(r : R_A\,a\,a').\,R_B\,a\,a'\,r\,(f a)\,(f' a')$.
  $\texttt{forallVariance}$ is the analogue of \texttt{arrowVariance}; its
  wrinkle is that at cov $\geq 2a$ the domain is needed at the full equivalence
  (\texttt{map4}), so the forward map can produce the relatedness proof the
  codomain fiber is indexed by.
\end{definition}

\begin{theorem}[The graded dependent $\Pi$]
  \label{thm:param-forall}
  \lean{LeTrocq.paramForall}
  \leanok
  \uses{def:forall-variance, def:param, def:weaken, lem:map4-subsingleton}
  $\texttt{paramForall}\;m\;n$ builds
  $\Param\;\cls{m}{n}\;(\forall a, B\,a)\;(\forall a', B'\,a')$ from a domain
  witness and a codomain \emph{family} of witnesses, each at the
  \texttt{forallVariance}-minimal class. Over \SortU\ (so $B$ may land in
  \PropSort), whence $\forall x, P\,x$ for a \PropSort-valued $P$ transfers.
  \texttt{R\_in\_map} uses the domain's completeness to \texttt{subst} the fiber
  and \Cref{lem:map4-subsingleton} to identify the two relatedness proofs.
\end{theorem}

\section{Grading: register strong, use minimal}
\label{sec:grading}

A witness is registered once at a \textbf{strong} class (typically the
equivalence $\cls{4}{4}$); each occurrence is used at the \textbf{minimal} class
it needs; a weakening map (\Cref{def:weaken}) bridges the two. The minimal
classes are computed \textbf{top-down in a single pass}: the demanded output
class flows down through the type's dependency structure, and each former's
grading table (\Cref{def:arrow-variance}, \Cref{def:forall-variance}) dictates
the class each part is built at --- no constraint graph, no fixpoint
(\Cref{chap:solver}). A leaf is built at its available class and \emph{weakened}
to the demand. One exception: a \emph{bound type variable} is pinned at
$\cls{4}{4}$ (the universe's inner class, \Cref{thm:param-type}) rather than the
join of its uses --- the top weakens to satisfy every use, and pinning it is
exactly what removes the fixpoint.

\section{Leaves: registered witnesses}
\label{sec:leaves}

The structural rules bottom out at \texttt{@[trocq]}-registered constants; the
translation never unfolds a definition (an unregistered head is an error).

\begin{definition}[Witness classification]
  \label{def:reg-kind}
  \lean{LeTrocq.parseEntry}
  \leanok
  \uses{def:param}
  Tagging a constant \texttt{@[trocq]} runs \texttt{parseEntry}, which reads the
  conclusion of its telescoped type into a \texttt{RegKind}: a \textbf{base}
  ($\Param\;m\;n\;A\;B$, closed $A,B$), a \textbf{relator}
  ($\forall \ldots, \Param\;m\;n\;(P\ldots)\,(P'\ldots)$ --- also how a
  \PropSort\ predicate or connective registers), a \textbf{type former} (a
  parametricity relation, conclusion a \SortU), or a \textbf{term primitive}
  (bare-relation conclusion). Witnesses are stored by name and re-created with
  fresh universe levels, so universe-polymorphic witnesses register correctly.
\end{definition}

\begin{definition}[The counterpart translation]
  \label{def:translate-term}
  \lean{LeTrocq.Translate.term}
  \leanok
  \uses{def:reg-kind}
  $\Cpt{\cdot}$ (\texttt{Translate.term}) rebuilds a term's $B$-side counterpart
  leaf by leaf: a registered head maps to its counterpart, a bound variable to
  $x'$, application/$\lambda$/$\Pi$/sort structurally, and a \texttt{Nat} numeral
  through its \texttt{succ}/\texttt{zero} normal form. It is the sole consumer's
  companion to $\Rel{\cdot}$ ($\Rel{t\,u} = \Rel{t}\,u\,\Cpt{u}\,\Rel{u}$).
\end{definition}

A parameterized type former $F$ crosses through its relator by the very same
abstraction theorem, with $\Rel{F}$ a $\Param$-valued witness and a final
weakening:
\[
  \Rel{F\,a_1 \ldots a_n} := \big(\Rel{F}\; a_1\,\Cpt{a_1}\,\Rel{a_1}\; \ldots\; a_n\,\Cpt{a_n}\,\Rel{a_n}\big)\ \text{weakened to the demand}.
\]

% ===========================================================================
\chapter{The demand-driven witness solver}
\label{chap:solver}

The relational translation $\Rel{\cdot}$ is realised as a single syntax-directed
pass that \textbf{produces the $\Param$ witness directly}, driven by a demanded
output class: no constraint graph, no least fixpoint. It has two
mutually-recursive halves --- a top-down \emph{type} judgment and a bottom-up
\emph{term} judgment.

\begin{definition}[The three judgments]
  \label{def:judgments}
  \uses{def:param, def:sem}
  \[
    \begin{array}{lll}
      \text{Type check} & \Gamma \vdash T \chk \cls{n}{m} \rw p & p : \Param\;n\;m\;T\;\Cpt{T} \quad(\text{top-down})\\
      \text{Type leaf}  & \Gamma \vdash T \syn \cls{n}{m} \rw p & (\text{atoms / type-vars: read the available class})\\
      \text{Term}       & \Gamma \vdash t \rw r                 & r : \Sem{\mathrm{typeof}\,t}\; t\; \Cpt{t}
    \end{array}
  \]
  A type variable enters the context as $A \mapsto (\Cpt{A} = A',\ a_R : \Param\;4\;4\;A\;A')$
  (pinned strong); a term variable as $x \mapsto (x',\ x_R : \Sem{T}\,x\,x')$.
  Both halves share one environment (\texttt{SEnv}); a type variable's
  \texttt{fvar} has a \SortU\ type, a term variable's does not, so the two never
  confuse the judgments.
\end{definition}

\begin{definition}[Subsumption --- the only mode switch]
  \label{def:subsumption}
  \uses{def:weaken, def:judgments}
  \[
    \irule{ \Gamma \vdash T \syn \cls{n'}{m'} \rw p \qquad \cls{n}{m} \leq \cls{n'}{m'} }
          { \Gamma \vdash T \chk \cls{n}{m} \rw \mathrm{weaken}\;p }
  \]
  Used to satisfy a check demand from a leaf's inferred available class.
\end{definition}

\section{The type judgment (check, top-down)}
\label{sec:type-rules}

\begin{definition}[The type check pass]
  \label{def:assemble}
  \lean{LeTrocq.Transfer.assemble}
  \leanok
  \uses{def:judgments, def:subsumption, thm:param-arrow, thm:param-forall, thm:param-type, thm:param-prop, def:mk-univ, def:reg-kind, def:arg-kind}
  $\texttt{assemble} : \mathsf{Expr} \to \mathsf{ParamClass} \to \mathsf{MetaM}\,\mathsf{Expr}$
  walks a type top-down, building its $\Param$ witness directly at the demanded
  class. The structural rules push the demand through the variance tables to the
  minimal part classes and build each node with its combinator:
  \[
    \irule{ }{ \Gamma \vdash \PropSort \chk \cls{n}{m} \rw \texttt{paramPropAt}\;n\;m }
    \qquad
    \irule{ \cls{n}{m} \leq \cls{2a}{2a} }
          { \Gamma \vdash \Type_w \chk \cls{n}{m} \rw \texttt{paramTypeAtInner}\;n\;m\;4\;4 }
  \]
  \[
    \irule{ \Gamma \vdash A \chk \texttt{arrowVariance}\cls{n}{m}.1 \rw p_a \qquad
            \Gamma \vdash B \chk \texttt{arrowVariance}\cls{n}{m}.2 \rw p_b }
          { \Gamma \vdash A \to B \chk \cls{n}{m} \rw \texttt{paramArrow}\;n\;m\;p_a\;p_b }
  \]
  \[
    \irule{ \Gamma \vdash D \chk \texttt{forallVariance}\cls{n}{m}.1 \rw p_D \qquad
            \Gamma, x{:}D \vdash B \chk \texttt{forallVariance}\cls{n}{m}.2 \rw p_b }
          { \Gamma \vdash \Pi x{:}D.\,B \chk \cls{n}{m} \rw \texttt{paramForall}\;n\;m\;p_D\;(\lambda x\,x'\,x_R.\,p_b) }
  \]
  where a $\Pi$ over a \emph{type} domain ($D = \Type_w$) offers $x$ at inner
  class $\cls{4}{4}$ and a $\Pi$ over a \emph{term} domain reads the two sides
  off $p_D$'s type. A leaf --- a registered atom, or a bound type variable
  offered at $\cls{4}{4}$ --- reads its available class and weakens
  (\Cref{def:subsumption}). A demand $\Type \chk \cls{n}{m}$ with
  $\cls{n}{m} \not\leq \cls{2a}{2a}$ \textbf{fails}: that is exactly ``needs
  univalence''.
\end{definition}

\begin{definition}[The inner class is pinned]
  \label{def:inner-class}
  \lean{LeTrocq.Transfer.innerClass}
  \leanok
  \uses{thm:param-type}
  Every universe combinator carries its bound type variable at inner class
  $\cls{4}{4}$: the top, independent of the capped outer class, weakening to
  satisfy every use. Pinning it is what lets a bound variable have a fixed class
  --- no fixpoint --- and it is axiom-free ($\texttt{paramIdAt}\;\cls{4}{4}$ is
  \texttt{paramRefl} weakened).
\end{definition}

\begin{definition}[The universe cap, enforced]
  \label{def:mk-univ}
  \lean{LeTrocq.Transfer.mkUniv}
  \leanok
  \uses{thm:param-type}
  $\texttt{mkUniv}$ emits the universe combinator at the demanded outer class,
  \textbf{failing} when it exceeds $\cls{2a}{2a}$. This is where feasibility of a
  $\Type$ transfer is decided; there is no separate sort table.
\end{definition}

\begin{definition}[Relator argument routing]
  \label{def:arg-kind}
  \lean{LeTrocq.Solver.relatorArgKinds}
  \leanok
  \uses{def:reg-kind}
  For an application $F\,a_1\ldots a_k$ of a registered relator, the abstraction
  theorem $\Rel{F\,\bar a} = \Rel{F}\,a_1\Cpt{a_1}\Rel{a_1}\ldots$ (then weaken to
  the demand) consumes each argument by its \emph{kind}, read off the relator's
  type by the shape of each triple's relatedness: \texttt{.type} (a $\Param$: a
  type argument, assembled at the relator's declared class), \texttt{.family}
  (relatedness telescopes to a $\Param$: a dependent type family, built like a
  $\Pi$ codomain), or \texttt{.term} (a bare relation: routed to the term
  judgment). A family reads its domain off its own binder type, so it need not
  sit adjacent to that domain.
\end{definition}

\section{The term judgment (abstraction theorem, bottom-up)}
\label{sec:term-rules}

\begin{definition}[The term pass]
  \label{def:assemble-term}
  \lean{LeTrocq.Transfer.assembleTerm}
  \leanok
  \uses{def:assemble, def:translate-term, def:param-of-iff, def:judgments}
  $\texttt{assembleTerm}$ is the abstraction theorem, bottoming at registered
  term primitives:
  \[
    \irule{ x \mapsto (\_, x_R) \in \Gamma }{ \Gamma \vdash x \rw x_R }
    \qquad
    \irule{ c\ \text{a registered primitive} \rw w }{ \Gamma \vdash c \rw w }
  \]
  \[
    \irule{ \Gamma \vdash u \rw u_R \qquad \Gamma \vdash v \rw v_R }
          { \Gamma \vdash u\,v \rw u_R\; v\; \Cpt{v}\; v_R }
  \]
  \[
    \irule{ \Gamma \vdash A \chk \cls{0}{0} \rw p_A \quad (\Sem{A} = p_A.R) \qquad \Gamma, x{:}A \vdash t \rw t_R }
          { \Gamma \vdash \lambda x{:}A.\,t \rw \lambda x\,x'\,x_R.\,t_R }
  \]
  \[
    \irule{ t : \PropSort \qquad \Gamma \vdash t \chk \cls{1}{1} \rw p }
          { \Gamma \vdash t \rw \mathsf{PLift.up}\,(\texttt{iffOfParam}\;p) }
    \qquad
    \irule{ \text{numeral } n }{ \Gamma \vdash n \rw (\text{expand to succ/zero}) }
  \]
  A term position consumes only $\Sem{A} = p_A.R$, which is grade-invariant, so
  the $\lambda$ rule checks its domain at the cheapest class $\cls{0}{0}$. A
  proposition is just a $\SortU\,0$ type, so it has no separate arm: its
  relatedness $\mathsf{PLift}\,(P \leftrightarrow P')$ is the $\cls{1}{1}$ witness
  the relator path builds, projected via \texttt{iffOfParam}.
\end{definition}

\section{Feasibility and entry points}
\label{sec:entry}

\begin{remark}[Feasibility falls out]
  \label{rem:feasibility}
  \uses{def:assemble, def:mk-univ, def:subsumption}
  A transfer is impossible exactly when a check gets stuck: a $\Type$ demand
  above $\cls{2a}{2a}$ (\Cref{def:mk-univ}), or a leaf whose available class
  cannot weaken to the demand. Both surface as ordinary rule failures, so there
  is no separate feasibility pass.
\end{remark}

\begin{definition}[The entry points]
  \label{def:entry}
  \lean{LeTrocq.Transfer.transfer, LeTrocq.Transfer.relate}
  \leanok
  \uses{def:assemble, def:assemble-term}
  \texttt{transfer} assembles $\Rel{T}$ directly at a root class in one pass
  (driving \texttt{transfer\%} at $\cls{4}{4}$ and the \texttt{trocq} tactic at
  the comap class $\cls{0}{1}$, which refines the goal by the backward map);
  \texttt{relate} runs the term judgment (driving \texttt{relate\%}).
  \texttt{translate\%} is $\Cpt{\cdot}$ (\Cref{def:translate-term}), unchanged.
\end{definition}
